\section*{Введение}
\addcontentsline{toc}{section}{Введение}

\textbf{Актуальность темы.} В условиях стремительного роста рынка электронной коммерции 
спрос на онлайн-покупку спортивных товаров неуклонно увеличивается. Традиционные 
офлайн-магазины не способны в полной мере удовлетворить потребности современного 
покупателя, ориентированного на удобный доступ к широкому ассортименту, 
персонализированные рекомендации и прозрачное управление заказами. Разработка 
специализированных веб-приложений интернет-магазинов, объединяющих функциональный 
каталог товаров, систему управления заказами и личные кабинеты пользователей, является 
необходимым условием для повышения конкурентоспособности компаний в сфере спортивной
розничной торговли.

\textbf{Степень разработанности проблемы.} Задача построения интернет-магазинов с 
полноценной системой управления заказами широко рассматривается в современной литературе 
по веб-разработке и электронной коммерции. Основы проектирования пользовательских личных 
кабинетов и систем обработки заказов заложены в работах по архитектуре клиент-серверных 
приложений и UX-проектированию. В области веб-ориентированных торговых платформ существенный 
вклад внесли исследователи в сфере баз данных, REST API и систем аутентификации пользователей. 
Несмотря на наличие готовых коммерческих решений (напр. Shopify, WooCommerce), разработка 
гибких специализированных приложений для нишевых рынков, таких как спортивные товары, с 
учётом специфики ассортимента и поведения целевой аудитории, остаётся актуальной и не 
до конца решённой практической задачей.

\textbf{Объект исследования} — процессы организации и управления онлайн-торговлей 
спортивными товарами, включая взаимодействие покупателей с интернет-магазином 
и обработку заказов.

\textbf{Предмет исследования} — программные средства и архитектурные решения для разработки 
веб-приложения интернет-магазина спортивных товаров, обеспечивающего автоматизацию управления 
заказами и функциональность личных кабинетов пользователей.

\textbf{Целью данной работы} является проектирование и реализация клиент-серверного веб-приложения
интернет-магазина спортивных товаров, обеспечивающего автоматизацию управления заказами, 
ведение каталога продукции и предоставление пользователям персональных личных кабинетов.
